\chapter{Mathematische Grundlagen}

\section{Einleitung}

Dieser Anhang stellt die mathematischen Werkzeuge bereit, auf denen die diskrete Formulierung der IWT basiert.

Im Gegensatz zu den kontinuierlichen Approximationen im Haupttext werden hier ausschließlich die fundamentalen diskreten Strukturen entwickelt.

\section{Diskrete Variationsrechnung}

\subsection{Das diskrete Funktional}

Ein diskretes Informationssystem wird beschrieben durch Knoten \( k = 1, \dots, N \) mit komplexen Informationswerten \( I_k^{(n)} \) zum Zeitschritt \( n \).

Das fundamentale diskrete Funktional lautet:

\[
\mathcal{S}_d[I_k^{(n)}] = \sum_{n=0}^{M-1} \sum_{k=1}^{N} \mathcal{F}_d\left(I_k^{(n)}, \Delta_t I_k^{(n)}, \Delta_s I_k^{(n)}\right) \cdot T \cdot V_k
\]

mit:
\begin{itemize}
\item \( \Delta_t I_k^{(n)} = I_k^{(n)} - I_k^{(n-1)} \) – zeitliche Differenz
\item \( \Delta_s I_k^{(n)} = \sum_{l \in \mathcal{N}(k)} w_{kl}(I_l^{(n)} - I_k^{(n)}) \) – räumliche Differenz
\item \( T \) – fundamentaler Zeitschrift
\item \( V_k \) – diskretes Volumenelement am Knoten \( k \)
\end{itemize}

\subsection{Diskrete Euler-Lagrange-Gleichungen}

Die Variation des Funktionals:

\[
\delta \mathcal{S}_d = 0
\]

führt zu den diskreten Euler-Lagrange-Gleichungen:

\[
\frac{\partial \mathcal{F}_d}{\partial I_k^{(n)}} - \Delta_t^+ \left( \frac{\partial \mathcal{F}_d}{\partial (\Delta_t I_k^{(n)})} \right) - \Delta_s \left( \frac{\partial \mathcal{F}_d}{\partial (\Delta_s I_k^{(n)})} \right) = 0
\]

mit:
\begin{itemize}
\item \( \Delta_t^+ f^{(n)} = \frac{f^{(n+1)} - f^{(n)}}{T} \) – diskrete Vorwärtsdifferenz
\item \( \Delta_s f_k = \sum_{l \in \mathcal{N}(k)} w_{kl}(f_l - f_k) \) – diskreter Laplace-Operator
\end{itemize}

\section{Diskrete Noether-Theoreme}

Eine diskrete Transformation:

\[
I_k^{(n)} \to I_k^{(n)} + \epsilon \cdot \xi_k^{(n)}
\]

ist eine Symmetrie, wenn das Funktional invariant bleibt.

Aus jeder solchen Symmetrie folgt eine diskrete Erhaltungsgröße:

\[
Q^{(n+1)} - Q^{(n)} = -\sum_k \Delta_s J_k^{(n)}
\]

mit:

\[
Q^{(n)} = \sum_k \left[ \xi_k^{(n)} \frac{\partial \mathcal{F}_d}{\partial (\Delta_t I_k^{(n)})} - K^{(n)} \right]
\]

\subsection{Wichtige Symmetrien und Erhaltungssätze}

\begin{itemize}
\item \textbf{Zeittranslation:} \( I_k^{(n)} \to I_k^{(n+1)} \) – Erhaltung der Energie \( E^{(n)} \).
\item \textbf{Raumtranslation:} \( I_k^{(n)} \to I_{k+\delta}^{(n)} \) – Erhaltung des Impulses \( \vec{P}^{(n)} \).
\item \textbf{Rotation:} \( I_k^{(n)} \to I_{R(k)}^{(n)} \) – Erhaltung des Drehimpulses \( \vec{L}^{(n)} \).
\item \textbf{Informations-Shift:} \( I_k^{(n)} \to I_k^{(n)} + c \) – Erhaltung der Gesamtinformation \( \sum |I|^2 \).
\end{itemize}

\section{Die diskrete Informationsmetrik}

Die diskrete Informationsmetrik zwischen Knoten \( k \) und \( l \) ist definiert als:

\[
g_{kl}^{(n)} = \frac{\partial^2 \mathcal{F}_d}{\partial (\Delta_s I_k^{(n)}) \partial (\Delta_s I_l^{(n)})}
\]

Alternativ über die fraktale Kopplungsmatrix:

\[
g_{kl}^{(n)} = \frac{K_{kl}^{(n)}}{\sqrt{K_{kk}^{(n)} K_{ll}^{(n)}}}
\]

mit:

\[
K_{kl}^{(n)} = \frac{1}{d_{kl}^{3-D}}
\]

\subsection{Eigenschaften}

\begin{itemize}
\item Symmetrie: \( g_{kl}^{(n)} = g_{lk}^{(n)} \)
\item Positive Definitheit: \( \sum_{k,l} v_k g_{kl}^{(n)} v_l \ge 0 \) für alle \( v_k \)
\item Normierung: \( g_{kk}^{(n)} = 1 \)
\item Lokalität: \( g_{kl}^{(n)} \to 0 \) für \( |k-l| \to \infty \)
\end{itemize}

\section{Fraktale Analysis diskreter Netzwerke}

\subsection{Fraktale Dimension}

Für eine selbstähnliche Struktur gilt:

\[
N(R) \propto R^D
\]

wobei \( N(R) \) die Anzahl der Knoten innerhalb einer Distanz \( R \) von einem Referenzknoten ist.

\subsection{Die Kopplungsmatrix}

Die Kopplungsmatrix des fraktalen Netzwerks ist:

\[
K_{kl} = \frac{1}{d_{kl}^{3-D}}
\]

wobei \( d_{kl} \) die fraktale Distanz im Indexraum ist.

\section{Diskrete Fourier-Analyse}

Für ein diskretes Feld \( I_k \) auf \( N \) regelmäßig angeordneten Knoten:

\[
\tilde{I}_q = \frac{1}{\sqrt{N}} \sum_{k=0}^{N-1} I_k e^{-2\pi i q k / N}
\]

Die Dispersionsrelation für das fraktale Netzwerk mit spektraler Dimension \( d_s = 2D/3 \):

\[
\omega \propto k^{D/3}
\]

\section{Zusammenfassung der mathematischen Struktur}

\begin{itemize}
\item Zustandsraum: Diskrete, komplexe Felder \( I_k^{(n)} \) auf Netzwerkknoten.
\item Dynamik: Diskretes Funktional \( \mathcal{S}_d[I_k^{(n)}] \) mit Variationsprinzip.
\item Bewegungsgleichungen: Diskrete Euler-Lagrange-Gleichungen.
\item Symmetrien: Diskrete Noether-Theoreme mit Erhaltungssätzen.
\item Geometrie: Diskrete Metrik \( g_{kl}^{(n)} = K_{kl}^{(n)} / \sqrt{K_{kk}^{(n)} K_{ll}^{(n)}} \).
\item Kopplung: Fraktale Kopplungsmatrix \( K_{kl}^{(n)} = 1/d_{kl}^{3-D} \).
\item Skalierung: Fraktale Dimension \( D \) charakterisiert das Netzwerk.
\end{itemize}